% !PRINTED 485
\weberpage{485}
%<<CONT
and if one adds the two equations (35):
\[
\begin{aligned}
&[(\xi+\eta)^3-3\,\xi\eta(\xi+\eta)]\left(1+\frac{1}{\xi^3\eta^3}\right)\\
&+\,2[(\xi+\eta)^2-2\,\xi\eta]\left(1-\frac{1}{\xi^2\eta^2}\right)=0.
\end{aligned}
\]

But by (34):
\[
\sqrt2\,\xi\eta=x^2,
\]
and one thus obtains an equation for $x^2$ when one eliminates $(\xi+\eta)$. By the ordinary elimination procedure this equation is found, when one sets
\[
z=\frac{x^2}{\sqrt2}+\frac{\sqrt2}{x^2}
\]
in the form
\weq{36}{z^6-9\,z^5+20\,z^4+6\,z^3-19\,z^2-17\,z-6=0}
But this still contains the factor relating to the determinant $-49$. For this factor one has [by formula (10)]:
\[
z=2+\sqrt7.
\]

The left side of (36) must therefore be divisible by
\[
z^2-4\,z-3
\]
and the carrying out of the division yields the factor valid for the determinant $-41$:
\weq{37}{\begin{aligned}
z^4-5\,z^3+3\,z^2+3\,z+2&=0, \\
z&=\dfrac{f(\sqrt{-41})}{\sqrt2}+\dfrac{\sqrt2}{f(\sqrt{-41})}.
\end{aligned}}

\begin{center}
\S\,131. \gesperrt{Calculation of class invariants from the irrational forms of the modular equations.}
\end{center}

In an extraordinarily simple way the irrational forms of the modular equations lead in many cases to the setting up of class equations.

1. In \S\,75 we have seen that, if $m\equiv -1\ (\mathrm{mod}\ 8)$, then between the two functions:
\[
\begin{aligned}
2A&=f(\omega)f(m\omega)+(-1)^{\frac{m+1}{8}}[f_1(\omega)f_1(m\omega)+f_2(\omega)f_2(m\omega)],\\
B&=\frac{2}{f_1(\omega)f_1(m\omega)}+\frac{2}{f_2(\omega)f_2(m\omega)}+\frac{(-1)^{\frac{m+1}{8}}2}{f(\omega)f(m\omega)}
\end{aligned}
\]
%CONT>>
